\documentclass[12pt]{article}
\usepackage[margin=1in]{geometry}
\usepackage{amsmath,amssymb,amsthm,amsfonts}
\usepackage{graphicx}
\usepackage{hyperref}
\usepackage{framed}
\hypersetup{colorlinks=true, linkcolor=blue, urlcolor=blue, citecolor=blue}

\theoremstyle{definition}
\newtheorem{definition}{Definition}[section]
\newtheorem{remark}{Remark}[section]
\theoremstyle{plain}
\newtheorem{theorem}{Theorem}[section]
\newtheorem{lemma}{Lemma}[section]
\newtheorem{corollary}{Corollary}[section]

\newenvironment{editorsnote}{%
  \begin{framed}
  \noindent\textbf{Editor's Note:}\em
}{\end{framed}}

\title{\textbf{Chapter 5: The Geometric Argument: Testability and its Limits}}
\author{Dmytro Surdu}
\date{}

\begin{document}

\maketitle

\begin{editorsnote}
This chapter presents the first major synthesis of our theory. Using the Measurement Chain model from Chapter 4 and the geometric tools from Chapter 3, this paper builds the complete information-theoretic case for our thesis. This is a revised and polished version of the original draft, now employing the more powerful metric-entropy-based arguments and a more precise definition of an admissible predicate.

It accomplishes three critical tasks:
\begin{enumerate}
    \item \textbf{It formalizes the witness.} The certificate for a property is defined as a `(coarse-label, fibre-index)` pair, which serves to pin down the hidden initial state of the system.
    \item \textbf{It proves the positive result (sufficiency).} The main theorem demonstrates that if a measurement model has positive geometric information loss ($\Delta^{\text{ent}}_\varepsilon > 0$) and satisfies a reasonable regularity condition, then tail properties can be certified by a short, polynomial-size $\mathbf{NP}$ witness.
    \item \textbf{It proves the negative result (necessity of loss).} It shows that if the model is perfectly lossless ($\Delta^{\text{ent}}_\varepsilon = 0$), the certification problem becomes provably intractable, placing it outside of $\mathbf{NP}$.
\end{enumerate}
This chapter provides the core "if and only if" argument, linking the geometric property of information loss directly to the logical property of testability.
\end{editorsnote}

\section*{Abstract}
We formalise and prove an ``information--loss $\Leftrightarrow$ testability'' theorem for deterministic measurement chains. We define information loss in purely geometric terms using the metric entropy of the system's state-space fibres. We show that a positive, uniformly bounded information loss—a property of non-invertible and regular (continuous, finite-to-one) models—is a sufficient condition for tail properties of the measurement stream to be certifiable in the complexity class $\mathbf{NP}$. Conversely, we show that a lossless channel leads to a provably intractable ($\Pi^0_2$-hard) certification problem. The argument is constructive, yielding an explicit finite witness that consists of a coarse system state and an index within the information-loss fibre.

\section{Introduction}

Complexity theory distinguishes problems that are merely \emph{testable} ($\mathbf{NP}$) from those that seem to require intractable proof (co-$\mathbf{NP}$ and harder). In metrology, we face an analogous gap: certain statements about a system can be certified from finite data, while others seem to require impossible knowledge of an infinite future. This paper proves that the bridge between these two worlds is **information loss**.

We will demonstrate that a deterministic measurement system permits a finite, polynomial-size certificate for a tail property if and only if the measurement channel is non-invertible. We formalize this non-invertibility as a positive, uniformly bounded geometric information loss, and show that the certificate is precisely a state-space coordinate: a coarse label plus an index locating the true micro-state within its collapsed fibre.

Our argument proceeds as follows. In Section \ref{sec:framework}, we define the necessary regularity conditions for a model and a predicate. In Section \ref{sec:main_results}, we define the state witness, prove our main sufficiency theorem (Theorem \ref{thm:sufficiency}), and then prove the corresponding hardness theorem for the lossless case (Theorem \ref{thm:hardness}).

\section{Framework}\label{sec:framework}
We adopt the Measurement Chain model defined in Chapter 4, with its components $Q, f, g$, and the resulting stream $Z_{1:\infty}$. We also use the definitions of the fibre $\mathcal{G}_y = (g \circ f)^{-1}(y)$ and geometric information loss $\Delta^{\text{ent}}_\varepsilon = \sup_{y \in \mathcal{Y}} \log_2 N_\varepsilon(\mathcal{G}_y)$ from Chapter 3.

\subsection{Model Regularity}
To ensure our results hold in the worst-case, we require a regularity condition on the measurement map $g \circ f$. This condition ensures the instrument's uncertainty structure is well-behaved and non-pathological.

\begin{definition}[Model Regularity (Finite-to-One on a Compact Domain)]
The measurement map $g \circ f$ satisfies \textbf{Model Regularity} if:
\begin{enumerate}
    \item The space of true quantities, $\mathcal{Q}$, is compact under the metric $d$.
    \item The map $g \circ f: \mathcal{Q} \to \mathcal{Y}$ is Lipschitz on $\mathcal Q$.
    \item The map is \textbf{finite-to-one}: there exists an integer $M$ such that for every $y \in \mathcal{Y}$, the fibre $\mathcal{G}_y$ has cardinality $|\mathcal{G}_y| \le M$.
\end{enumerate}
\end{definition}

\begin{lemma}[Consequence of Model Regularity]
If a model satisfies the regularity condition, then its geometric information loss $\Delta^{\text{ent}}_\varepsilon$ is uniformly bounded. For all sufficiently small $\varepsilon$, we have $N_\varepsilon(\mathcal{G}_y) = |\mathcal{G}_y| \le M$ for every fibre $y$. Therefore,
\[
\Delta^{\text{ent}}_\varepsilon \le \log_2 M
\]
The information loss is bounded by a constant determined by the model's structure.
\end{lemma}

\subsection{Admissible Predicates}
Our results apply to tail predicates that are "well-behaved" enough to be decided by a finite simulation and robust to small errors. This captures the requirement that the predicate must eliminate all residual randomness in finite time.

\begin{definition}[Admissible Tail Predicate]
A predicate $P_{\mathrm{tail}}$ on infinite streams $Z_{1:\infty}$ is **admissible** if there exist a polynomial $T(n)$ and a deterministic Boolean function $\Phi$ such that for any true state $q \in \mathcal{Q}$ and any $\varepsilon$-approximation $q^*$ with $d(q, q^*) \le \varepsilon$:
\begin{enumerate}
    \item \textbf{Finite-Window Determination:} The truth of the predicate is determined by a finite portion of the tail stream:
    \[
     P_{\mathrm{tail}}(Z_{1:\infty}) = \Phi(Z_{n+1}, \dots, Z_{n+T(n)})
    \]
    \item \textbf{$\varepsilon$-Robustness:} Small errors in the initial state do not flip the decision. Let $Z$ be the stream from $q$ and $Z^*$ be the stream from $q^*$. Then:
    \[
    \Phi(Z_{n+1}, \dots, Z_{n+T(n)}) = \Phi(Z^*_{n+1}, \dots, Z^*_{n+T(n)})
    \]
\end{enumerate}
\end{definition}

\section{The Main Results: Testability and its Limits}\label{sec:main_results}

\subsection{Sufficiency: Information Loss Enables NP-Certification}

\begin{definition}[State Witness]
Fix a precision $\varepsilon$. For any true state $q \in \mathcal{Q}$, its \textbf{state witness} is the pair $w(q) = (y, r)$, where:
\begin{itemize}
    \item $y = g(f(q))$ is the coarse instrument state (the fibre label).
    \item $r$ is the index of the $\varepsilon$-ball containing $q$ within a pre-defined, canonical cover of the fibre $\mathcal{G}_y$.
\end{itemize}
\end{definition}

\begin{lemma}[Uniform Witness‐Length Bound]
Under Model Regularity (compact $\mathcal Q$, Lipschitz $g\circ f$, finite‐to‐one with bound $M$), for any precision $\varepsilon>0$ the state witness
\[
w(q)\;=\;\bigl(y,\;r\bigr)
\]
where $y=g(f(q))$ and $r$ is the index of the $\varepsilon$‐ball containing $q$ in a fixed canonical cover of the fibre $\mathcal G_y$, has bit‐length
\[
|w(q)|\;\le\;\bigl\lceil\log_2 N_{\mathcal Y}(\varepsilon)\bigr\rceil\;+\;\lceil\log_2 M\rceil
\;=\;O\bigl(\log(1/\varepsilon)\bigr)\;+\;\log_2 M.
\]
Here $N_{\mathcal Y}(\varepsilon)$ is the size of a minimal $\varepsilon$‐cover of $\mathcal Y=g(f(\mathcal Q))$.
\end{lemma}

\begin{proof}
\noindent\textbf{1. Total boundedness of \(\mathcal Y\).}  
Since \(\mathcal Q\) is compact and \(g\circ f\) is Lipschitz, its image 
\(\mathcal Y = (g\circ f)(\mathcal Q)\) is also compact and hence totally bounded.  
By definition of total boundedness, for each \(\varepsilon>0\) there exists a finite set of “centers” 
\(\{y_1,\dots,y_{N_{\mathcal Y}(\varepsilon)}\}\subset\mathcal Y\) such that
\[
\mathcal Y \;\subset\;\bigcup_{j=1}^{N_{\mathcal Y}(\varepsilon)} B\bigl(y_j,\varepsilon\bigr).
\]
We fix once and for all this canonical \(\varepsilon\)‐cover and agree on a public indexing \(j=1,\dots,N_{\mathcal Y}(\varepsilon)\).

\medskip\noindent\textbf{2. Encoding the coarse label \(y\).}  
Given \(q\in\mathcal Q\), let \(y = g(f(q))\).  By the cover property, there is at least one index \(j\) with 
\[
d\bigl(y,y_j\bigr)\;\le\;\varepsilon.
\]
Choosing one (e.g.\ the smallest such \(j\)) identifies the \(\varepsilon\)‐ball containing \(y\).  
Encoding this index \(j\in\{1,\dots,N_{\mathcal Y}(\varepsilon)\}\) in binary requires
\[
\lceil \log_2 N_{\mathcal Y}(\varepsilon)\rceil
\]
bits.  By standard estimates for metric‐entropy, \(N_{\mathcal Y}(\varepsilon)=O(\varepsilon^{-d})\) in many spaces, so
\(\log_2 N_{\mathcal Y}(\varepsilon)=O(\log(1/\varepsilon))\).

\medskip\noindent\textbf{3. Encoding the fibre index \(r\).}  
Under the finite‐to‐one assumption, each fibre \(\mathcal G_y\) has size at most \(M\).  
We fix a public ordering of its elements and let \(r\in\{1,\dots,|\mathcal G_y|\}\) denote the position of the true state \(q\).  
Encoding \(r\) takes at most
\[
\lceil \log_2 M\rceil
\]
bits, which is a constant independent of \(\varepsilon\).

\medskip\noindent\textbf{4. Total bit‐length.}  
Concatenating the two binary codes for \(j\) and \(r\) yields the full witness \(w(q)\).  Hence
\[
|w(q)| \;\le\; \lceil \log_2 N_{\mathcal Y}(\varepsilon)\rceil \;+\;\lceil \log_2 M\rceil
\;=\; O\bigl(\log(1/\varepsilon)\bigr)\;+\;\log_2 M,
\]
as claimed.
\end{proof}


\begin{theorem}[Sufficiency]\label{thm:sufficiency}
Under Model Regularity, fix a precision $\varepsilon$ such that the geometric information loss $\Delta^{\text{ent}}_\varepsilon > 0$ (i.e., $M > 1$). Let $P_{\mathrm{tail}}$ be any admissible tail predicate. Then the language of "good heads" $L_P$ is in the complexity class $\mathbf{NP}$.
\end{theorem}

\begin{proof}
Our NP-certificate is precisely the state witness $w(q)=(y,r)$ corresponding to the true initial state $q$. Its length is polynomially bounded by the previous lemma. The verifier, on input head $z_{1:n}$ and witness $(y,r)$, performs the following polynomial-time algorithm:
\begin{enumerate}
  \item Reconstruct the center $q^\star$ of the $r$-th ball in the public cover of fibre $\mathcal{G}_y$. This $q^\star$ is an $\varepsilon$-approximation of the true state $q$.
  \item \textbf{Check Consistency:} Verify that simulating the stream generator starting from $q^\star$ for $n$ steps reproduces exactly the observed head $z_{1:n}$. If not, reject.
  \item \textbf{Verify Predicate:} Simulate the stream generator, starting from $q^\star$, for the next $T(n)$ steps (the horizon from the admissibility definition) to get the tail segment $z^\star_{n+1:n+T(n)}$.
  \item Evaluate the Boolean function $\Phi(z^\star_{n+1:n+T(n)})$. Accept if it returns true, reject otherwise.
\end{enumerate}
Because the predicate is $\varepsilon$-robust, the conclusion drawn from the simulation of $q^\star$ is the same as the one for the true state $q$. The verifier is deterministic, runs in polynomial time, and correctly accepts all true instances. Therefore, the language $L_P$ is in $\mathbf{NP}$.
\end{proof}

\subsection{Necessity: Lossless Channels are Intractable}
We now show that information loss is not just sufficient, but necessary. If the channel is lossless, the problem becomes provably intractable.

\begin{theorem}[Hardness]\label{thm:hardness}
Let the measurement map $g \circ f$ be a continuous bijection (i.e., it is invertible). This implies the channel is lossless, with $\Delta^{\text{ent}}_\varepsilon = 0$ for all $\varepsilon > 0$. For any non-trivial admissible tail predicate $P_{\mathrm{tail}}$, the language of "good heads" $L_P$ is $\Pi^0_2$-hard.
\end{theorem}
\begin{proof}[Sketch]
The full proof is detailed in the Appendix. The core idea is a reduction from the known $\Pi^0_2$-hard problem `TOTAL`. We construct a specific stream generator $h_e$ whose behavior depends on a Turing machine index $e$, and a predicate $P_e$ that is true if and only if $M_e$ halts on all inputs.

The key step is this: since $g \circ f$ is a bijection, the observed head $z_{1:n}$ uniquely determines the entire infinite stream’s initial state $q$. There is no ambiguity for a witness to resolve. The verifier is therefore faced with the full, unresolvable task of verifying the universal quantification 'for all inputs $x$...' required by the `TOTAL` reduction. This task is known to be $\Pi^0_2$-hard. Since we can reduce any instance of `TOTAL` to an instance of our problem, our problem must be at least as hard. It therefore cannot be in $\mathbf{NP}$ (unless the Polynomial Hierarchy collapses).
\end{proof}

\section{Conclusion}

The geometric properties of a measurement map determine its certificability. A non-invertible map that satisfies Model Regularity induces a positive, bounded geometric information loss. This loss creates the structure necessary for a short, finite witness to exist, placing the certification problem in $\mathbf{NP}$. Conversely, a perfectly lossless (invertible) map offers no such structure, leaving the certification problem in the intractable realm of the higher arithmetic hierarchy. Information loss is therefore the fundamental enabler of deterministic testability.

\end{document}